% Part: computability
% Chapter: computability-theory
% Section: complement-ce

\documentclass[../../../include/open-logic-section]{subfiles}

\begin{document}

% SEGMENT OLP-0242-S01

\olfileid{cmp}{thy}{cmp}
\olsection{கணிக்கத்தக்கவாறு எண்ணிடத்தக்க கணங்கள் நிரப்பியின்கீழ் மூடியவை அல்ல}

$A$ கணிக்கத்தக்கவாறு எண்ணிடத்தக்கது எனக் கொள்க. $A$ இன் நிரப்பி,
$\Complement{A} = \Nat \setminus A$, எப்போதும் கணிக்கத்தக்கவாறு
எண்ணிடத்தக்கதாகவும் இருக்குமா? விடை ``இல்லை'' என்று பின்வரும்
தேற்றமும் அதன் முடிவுரையும் காட்டுகின்றன.

\begin{thm}
\ollabel{thm:ce-comp}
$A$ இயல் எண்களின் ஏதாவது ஒரு கணம் ஆகட்டும். $A$ உம்
$\Complement{A}$ உம் கணிக்கத்தக்கவாறு எண்ணிடத்தக்கவை எனில், எனில்
மட்டுமே, $A$ கணிக்கத்தக்கது.
\end{thm}

\begin{proof}
முன்னோக்குத் திசை எளிது: $A$ கணிக்கத்தக்கது எனில்
$\Complement{A}$ உம் கணிக்கத்தக்கது
($\Char{A} = 1 \tsub \Char{\Complement{A}}$); எனவே இரண்டுமே
கணிக்கத்தக்கவாறு எண்ணிடத்தக்கவை.

மறுதிசையில், $A$ உம்~$\Complement{A}$ உம் கணிக்கத்தக்கவாறு
எண்ணிடத்தக்கவை எனக் கொள்க. $A$ ஆனது~$\cfind{d}$ இன்
வரையறைப்புலமாகவும், $\Complement{A}$ ஆனது~$\cfind{e}$ இன்
வரையறைப்புலமாகவும் இருக்கட்டும். $h$ ஐப் பின்வருமாறு வரையறுக்க:
\[
h(x) = \umin{s}{(T(d,x,s) \lor T(e,x,s))}.
\]
வேறுவிதமாகக் கூறினால், உள்ளீடு~$x$ இல், $\cfind{d}$ இன் நிற்கும்
கணிப்பு அல்லது $\cfind{e}$ இன் நிற்கும் கணிப்பு ஒன்றை~$h$ தேடுகிறது.
இப்போது $x \in A$ எனில் முதல் நேர்வில் அது வெற்றிபெறும்; $x \in
\Complement{A}$ எனில் இரண்டாவது நேர்வில் வெற்றிபெறும். ஆகவே $h$
ஒரு முழுக் கணிக்கத்தக்க சார்பு. ஆனால் இப்போது ஒவ்வொரு~$x$ க்கும்,
$T(e, x, h(x))$, அதாவது $\cfind{e}$ வரையறுக்கப்படுவது எனில், எனில்
மட்டுமே, $x \in A$. $T(e, x, h(x))$ ஒரு கணிக்கத்தக்க உறவாக இருப்பதால்,
$A$ கணிக்கத்தக்கது.
\end{proof}

\begin{explain}
முறைசாராத கணிப்பியல் மொழியில் நடப்பதைப் புரிந்துகொள்வது எளிது:
$A$ ஐத் தீர்மானிக்க, உள்ளீடு $x$ இல் $\cfind{e}$ மற்றும்
$\cfind{f}$ ஆகியவற்றின் நிற்கும் கணிப்புகளைத் தேடுக. அவற்றில் ஒன்று
கட்டாயம் நிற்கும்; அது $\cfind{e}$ எனில் $x$ ஆனது~$A$ இல் உள்ளது;
இல்லையெனில் $x$ ஆனது~$\Complement{A}$ இல் உள்ளது.
\end{explain}

\begin{cor}
\ollabel{cor:comp-k}
$\Complement{K_0}$ கணிக்கத்தக்கவாறு எண்ணிடத்தக்கது அல்ல.
\end{cor}

\begin{proof}
$K_0$ கணிக்கத்தக்கவாறு எண்ணிடத்தக்கது; ஆனால் கணிக்கத்தக்கது அல்ல
என்று அறிவோம். $\Complement{K_0}$ கணிக்கத்தக்கவாறு எண்ணிடத்தக்கதாக
இருந்தால், \olref{thm:ce-comp} ஆல் $K_0$ கணிக்கத்தக்கதாக இருக்கும்;
இது \olref[ncp]{thm:K-0} உடன் முரண்படுகிறது.
\end{proof}

\end{document}
